\section{Computations for the rank-five family}\label{sec:ppt3-app-family}

This appendix contains the proofs of the statements of Section~\ref{sec:ppt3-family}, whose notation it uses.

\begin{proof}[Proof of Theorem~\ref{thm:ppt3-family}, PPT and rank claims]
Cyclic permutations of the qubits leave $\rho_t$ invariant, so
its three partial transposes are unitarily equivalent. The
unnormalised numerator of $\rho_t^{T_A}$ commutes with the global parity $Z^{\otimes3}$, so it splits into even and
odd parity blocks. With $\beta=3+\omega$ and $\gamma=t+3\omega^2$ ($|\beta|^2=7$, $|\gamma|^2=t^2-3t+9$), in the respective basis orders
$(000,011,101,110)$ and $(001,010,100,111)$ these blocks of $7(t+3)\rho_t^{T_A}$ are
\[
\text{even: }\begin{pmatrix}
t & -t & \beta & \bar\beta\\
-t & t+3 & 0 & 0\\
\bar\beta & 0 & t+3 & \gamma\\
\beta & 0 & \bar\gamma & t+3
\end{pmatrix},\qquad
\text{odd: }\begin{pmatrix}
4 & \beta & -t & \bar\gamma\\
\bar\beta & 4 & -t & \gamma\\
-t & -t & 4 & 0\\
\gamma & \bar\gamma & 0 & 3t
\end{pmatrix},
\]
and their leading principal minors are
\begin{equation}\label{eq:ppt3-minors}
\begin{aligned}
\text{even: }&
t,\quad 3t,\quad (t+3)(3t-7),\quad
3(8t^2-15t-36),\\
\text{odd: }&
4,\quad 9,\quad 3(12-t^2),\quad
-3(t-3)(3t^2+4t-48).
\end{aligned}
\end{equation}
The even determinant is nonnegative exactly when $t\ge t_-$.
For $3\le t<2\sqrt3$, the first three leading minors
of each block are positive, while the odd determinant is
nonnegative exactly when $t\le t_+$. Since
$3<t_-<10/3<t_+<2\sqrt3$, the leading $3\times3$ corners are positive definite throughout $[t_-,t_+]$.
Each block is positive semidefinite exactly when its scalar Schur complement, or equivalently its determinant, is
nonnegative. This gives PPT precisely on $[t_-,t_+]$.
Outside this interval one of the displayed determinants
is negative. At either endpoint exactly one block has
rank three (its leading $3\times3$ minor stays positive);
the other has rank four. The stated profiles follow.
The GME claim is established in Section~\ref{sec:ppt3-witness}, after constructing
a witness on the range.
\end{proof}

\begin{proof}[Proof of Proposition~\ref{prop:ppt3-witness}]
\emph{(a).} The product condition reduces to four linear equations. Write a general vector of $R_s$ as
$B_s(U,V,W,E,O)^T$. It is a product with the first-qubit ray
$(1,z)^T$ precisely when the two $BC$ rows are proportional
with ratio $z$. In the computational basis this is
\[
\begin{aligned}
-zU+W+O&=0,&
-sU-zW+\omega E-z\omega^2O&=0,\\
-sU-zW+\omega^2E-z\omega O&=0,&
szU+V-zE&=0.
\end{aligned}
\]
Subtracting the middle equations and using
$\omega+\omega^2=-1$ gives
\[
 E=-zO,\qquad W=zU-O,\qquad
 (s+z^2)U=2zO,\qquad V=-szU-z^2O.
\]
For every finite $z$, the solutions form a one-dimensional
space, because $s+z^2$ and $2z$ never vanish simultaneously
when $s>0$. Setting $U=2z$, $O=s+z^2$ gives exactly
$q_s(z)$ in \eqref{eq:ppt3-curve}; it remains nonzero at
$z=0$ and $z^2=-s$. At the projective point
$(0,1)^T$ (i.e.\ $z=\infty$), the upper $BC$ row must
vanish, which forces $U=W=E=O=0$ and leaves only $V$,
the ray of $v$. This proves the full cut-product classification.

For Appendix~\ref{sec:ppt3-kernel}, $R^\perp$ is spanned by the three symmetry-sector vectors omitted from $R$, namely
$\bar e$, $\bar o$ and $3\ket{000}+\ket{011}+\ket{101}+\ket{110}$: the $3\times8$ matrix $K$ whose rows are their bras,
\[
K_0=[0,0,0,1,0,\omega,\omega^2,0],\quad
K_1=[0,\omega^2,\omega,0,1,0,0,0],\quad
K_2=[3,0,0,1,0,1,1,0]
\]
has rank three and $KB=0$, so $R=\ker K$.

\emph{(b).} A rescaling removes the parameter $s$: with
$D_s=\operatorname{diag}(1,s^{3/2},s^{1/2},s,s^{1/2})$,
\[
 q_s(\sqrt s\,z)=\sqrt s\,D_sq_1(z),\qquad
 H_s=s^3D_s^{-1}H_5D_s^{-1},
\]
so $q_s(\sqrt s z)\dg H_sq_s(\sqrt s z)=s^4\,q_1(z)\dg H_5q_1(z)$, and the right side of \eqref{eq:ppt3-sos} at
$(s,\sqrt s z)$ is $s^4$ times its value at $(1,z)$. It suffices to take $s=1$, which is \eqref{eq:ppt3-h5}. Put $d=\operatorname{Re}(z^2)=x^2-y^2$, so
$r+d=2x^2$. The squared moduli of the entries of $q_1(z)$ are
\[
4r,\quad r^2(r^2+9+6d),\quad r^2+1-2d,\quad r(r^2+1+2d),\quad r^2+1+2d,
\]
and $H_5$ is diagonal, so
\begin{align*}
q_1\dg H_5q_1&=-48r+r^2(r^2+9+6d)+11(r^2+1-2d)-6r(r^2+1+2d)+14(r^2+1+2d)\\
&=r^4-6r^3+34r^2-54r+25+6d(r-1)^2 .
\end{align*}
Since $(r-1)^2(r^2-4r+25)=r^4-6r^3+34r^2-54r+25$ and $12x^2=6(r+d)$, this equals
$(r-1)^2\bigl[(r-5)^2+12x^2\bigr]$. This proves \eqref{eq:ppt3-sos}.
The cyclic permutation of qubits fixes $u_s,v,w$ and
acts on $e,o$ by phases, which commute with the diagonal $H_s$, so the form controls the other cuts as well.
The point at infinity contributes $H_{vv}=1$. Finally, every positive pure-state decomposition of an operator
supported on $R_s$ uses vectors in $R_s$: its expectation
on $R_s^\perp$ vanishes, and each positive summand must
vanish there separately. Thus the nonnegative curve
values also control mixtures across different cuts.
\end{proof}

\begin{proof}[Proof of Corollary~\ref{cor:ppt3-chirality}]
\emph{(a).} The five nonzero eigenvalues of $\rho_t$ are, up to the common
factor $1/[7(t+3)]$, $4t,3t,9,9,3$, on the respective
rays $(u,v,w,e,o)$. Since $t>3$ throughout
\eqref{eq:ppt3-interval}, the two real eigenrays
$u$ and $v=\ket{111}$ are each nondegenerate.
Because $u$ and $v$ are real, $\overline{\rho_t}$ has the same
nondegenerate eigenrays, so a product unitary mapping $\rho_t$ to
its complex conjugate must fix the ray of $v=\ket{111}$. Each
factor then fixes the ray of $\ket1$, hence also the orthogonal
ray of $\ket0$, and is diagonal. Write its relative local phases as
$(\theta_A,\theta_B,\theta_C)$. Fixing the $u$ ray forces
\[
\theta_A+\theta_B\equiv\theta_A+\theta_C
\equiv\theta_B+\theta_C\equiv0\pmod{2\pi},
\]
so all three phases are $0$ or all are $\pi$.
Up to a scalar, the only candidates are $I$ and
$Z\otimes Z\otimes Z$, both of which fix $\rho_t$.
But, in the computational basis,
\[
  \langle011|\rho_t|101\rangle
       =\frac{t+3\omega^2}{7(t+3)}
       \notin\mathbb R .
\]
Thus $\rho_t\ne\overline{\rho_t}$, and neither candidate
implements complex conjugation.

\emph{(b).} By Proposition~\ref{prop:ppt3-witness}(a) the $A|BC$-product rays of $R_s$ are $[v]$ and
$[B_sq_s(z)]=[(\ket0+z\ket1)\otimes\ket{\phi_s(z)}]$, $z\in\mathbb C$, and substituting \eqref{eq:ppt3-curve} shows that the $2\times2$
coefficient matrix of $\ket{\phi_s(z)}$ has determinant $-3(z^2+s)^2$. So $\ket{\phi_s(z)}$ is a $B|C$ product exactly when
$z=\pm i\sqrt s$, where $B_sq_s(\pm i\sqrt s)=\pm2i\sqrt s\,(\ket0\pm i\sqrt s\ket1)^{\otimes3}$. Hence $R_s$ contains exactly three fully
product rays, $P_\infty=[\ket1^{\otimes3}]$ and $P_\pm=[(\ket0\pm i\sqrt s\ket1)^{\otimes3}]$; they span $\operatorname{span}(u_s,v,w)$, because
$(\ket0\pm i\sqrt s\ket1)^{\otimes3}=u_s\pm i\sqrt s\,(w-sv)$. The conjugate $\overline{R_s}$ contains the conjugate rays, the same three with $P_+$ and $P_-$ exchanged. An invertible
local filter $F$ that maps $R_s$ onto $R_s$ or onto $\overline{R_s}$ permutes the three rays. Each ray is $[\chi^{\otimes3}]$, so $F_A$,
$F_B$ and $F_C$ all induce the same permutation of the three distinct points $[\chi]\in\mathbb{CP}^1$, and since a M\"obius map is determined by its action on three
points, $F=cF_0^{\otimes3}$. Such an $F$ commutes with the cyclic shift $\Pi$, so it maps the $\omega^2$-eigenspace of $\Pi$ into itself. That eigenspace is
$\operatorname{span}(e,o)\subset R_s$, whereas $\overline{R_s}$ meets it in $0$: $\overline{R_s}\ni\bar e,\bar o$ lies in the $\omega$-eigenspace and the other three
generators of $\overline{R_s}$ are $\Pi$-invariant. So $F(R_s)\ne\overline{R_s}$. If a state $\rho$ with range $R_s$ satisfied $\rho=F\bar\rho F\dg$ (or the same with $Y^{\otimes3}\bar\rho Y^{\otimes3}$, a local-unitary image of $\bar\rho$), then
there would be an invertible local filter $F'$ ($F^{-1}$, respectively $Y^{\otimes3\dagger}F^{-1}$) with $F'(R_s)=\overline{R_s}$, which the argument excludes. For the stabiliser, each of the six permutations of the three points is
induced by a unique M\"obius map, and $F_\pi^{\otimes3}$ permutes $P_\infty,P_\pm$, hence preserves $\operatorname{span}(u_s,v,w)$, and as it commutes with $\Pi$ it preserves
the $\omega^2$-eigenspace $\operatorname{span}(e,o)$ of $\Pi$ (which is all of that eigenspace in $(\mathbb C^2)^{\otimes3}$). So the stabiliser is $\{[F_\pi]^{\otimes3}\}\cong S_3$.

\end{proof}

\begin{proof}[Proof of Corollary~\ref{cor:ppt3-extreme6}]
Write the numerator of $\sigma_6$ as $a\,uu\dg+b\,vv\dg+3\,ww\dg+3\,ee\dg+oo\dg$; it is the numerator of
$\rho_a$ with the weight $3a$ of $vv\dg$ replaced by $b$. The vector $v=\ket{111}$ enters $\sigma_6^{T_A}$ only through the
$(111,111)$ entry, the bottom diagonal entry of the odd block, and equals $b$ there. Hence the even block, and the three leading
principal minors of the odd block, are those of $\rho_a^{T_A}$ in \eqref{eq:ppt3-minors}: at $a=t_-$ the even minors
$a,\ 3a,\ (a+3)(3a-7)$ are positive and, since $8a^2-15a-36=0$, the even determinant vanishes, so the even block is positive
semidefinite of rank three. The odd minors $4,\ 9,\ 3(12-a^2)$ are positive ($t_-<2\sqrt3$), and the odd determinant is
affine in $b$,
\[
 \det(\text{odd block})=3(12-a^2)\,b+3(5a^2+24a-144),
\]
which equals $-3(t-3)(3t^2+4t-48)$ of \eqref{eq:ppt3-minors} at $a=t$, $b=3t$. It vanishes at
$b=(5a^2+24a-144)/(a^2-12)$, and this equals $8(9-a)/5$ modulo $8a^2-15a-36=0$, since
$5(5a^2+24a-144)-8(9-a)(a^2-12)=8a^3-47a^2+24a+144$ is divisible by $8a^2-15a-36$
(the quotient is $a-4$ and the remainder is $0$).
Hence at $b=8(9-a)/5$ the odd block is positive semidefinite of rank three as well. Thus each partial transpose of $\sigma_6$ has
rank six.
The supported witness has expectation
\[
\tr(H_5X_6)
=\frac{b-12a+29}{4a+b+21}
=\frac{217-68a}{5(4a+b+21)}<0,
\]
because $a>13/4>217/68$. This proves GME by
Proposition~\ref{prop:ppt3-witness}.
For extremality, $\sigma_6$ is extreme in the set of fully-PPT states exactly when the only Hermitian $D$ with range in $\operatorname{ran}\sigma_6$
and $D^{T_c}k=0$ for every kernel vector $k$ of $\sigma_6^{T_c}$, at each cut $c$, is a multiple of $\sigma_6$ (the supported-face
criterion of~\cite{GIM}). Write $D=BX_DB\dg$ with $X_D$ Hermitian. The conditions
$D^{T_c}k=0$ are a real-linear system on the $25$-dimensional real space of such $X_D$. To compute its rank exactly, complexify:
let $X_D$ range over all complex $5\times5$ matrices (25 complex unknowns) and impose both $D^{T_c}k=0$ and $k\dg D^{T_c}=0$ for the
kernel vectors $k$, at each cut. The two sets are adjoint to each other for Hermitian $D$, so the solution space is closed under
$X_D\mapsto X_D\dg$, and is the complexification of the real solution space on Hermitian $X_D$ (split $X_D$ into its
Hermitian and anti-Hermitian parts). Real nullity therefore equals complex nullity. The complexified system has rank $24$ over
$\mathbb Q(\sqrt{17},i\sqrt3)$, which contains the entries of $\sigma_6$ and of the kernel vectors, so its nullspace, and hence the
real nullspace, is exactly the scalar direction of \eqref{eq:ppt3-sigma6}.
\end{proof}
